% 06_tablas/matriz_jurisprudencial.tex
% Tabla 3. Matriz jurisprudencial — sentencias y documentos de la Corte Constitucional del Ecuador
% Capítulo II, sección 2.5.8 — Jurisprudencia de la Corte Constitucional
% Usa longtable para permitir extensión en varias páginas

\begingroup
\small
\begin{longtable}{|p{0.03\textwidth}|p{0.21\textwidth}|p{0.12\textwidth}|p{0.30\textwidth}|p{0.23\textwidth}|}

\caption{Matriz de jurisprudencia constitucional sobre derechos de la naturaleza, minería y derecho ambiental}
\label{tab:jurisprudencial} \\

\hline
\textbf{N.°} & \textbf{Resolución / Documento} & \textbf{Tipo} & \textbf{Holdings o contenido principal} & \textbf{Relevancia para la investigación} \\
\hline
\endfirsthead

\multicolumn{5}{l}{\small\textit{(continuación de la Tabla \ref*{tab:jurisprudencial})}} \\[2pt]
\hline
\textbf{N.°} & \textbf{Resolución / Documento} & \textbf{Tipo} & \textbf{Holdings o contenido principal} & \textbf{Relevancia para la investigación} \\
\hline
\endhead

\hline
\endfoot

1 &
Sentencia No.\ 1149-19-JP/21. Caso Bosque Protector Los Cedros. Corte Constitucional del Ecuador. Noviembre de 2021 &
Sentencia vinculante. Precedente constitucional &
(a) Estándares de la consulta ambiental (art.\ 398 CRE) en concesiones mineras. (b) Tres elementos del principio de precaución en el contexto extractivo. (c) \textit{In dubio pro natura} como criterio hermenéutico obligatorio. (d) Obligación de realizar la consulta antes de aprobar el EIA. &
Precedente más relevante para la tesis. Condiciona el licenciamiento ambiental, la consulta y la evaluación de compatibilidad de la apertura minera privada con el derecho ambiental constitucional. \\
\hline

2 &
Sentencia No.\ 32-17-IN/21. Acción de inconstitucionalidad de los arts.\ 86 y 136 del Reglamento Ambiental para Actividades Mineras. Corte Constitucional del Ecuador. 2021 &
Sentencia vinculante &
Principio de reserva de ley orgánica como límite a la regulación reglamentaria de materias que inciden en derechos constitucionales en el ámbito minero. &
Determina qué materias mineras y ambientales requieren ley orgánica y cuáles pueden regularse por reglamento. Relevante para el análisis de la técnica legislativa de las reformas. \\
\hline

3 &
Sentencia No.\ 22-18-IN/21. Acción de inconstitucionalidad de los arts.\ 462 y 463 del RCOAM. Corte Constitucional del Ecuador. 2021 &
Sentencia vinculante &
(a) Derechos de la naturaleza plenamente operativos en ecosistemas frágiles (manglares). (b) Distinción constitucional entre consulta ambiental (art.\ 398 CRE) y CPLI (art.\ 57.7 CRE): son mecanismos independientes. (c) Declaratoria de inconstitucionalidad de arts.\ 462--463 RCOAM por regular la CPLI mediante reglamento en violación de la reserva de ley orgánica. &
Tres holdings de relevancia sistémica. Fundamental para la delimitación entre consulta ambiental y CPLI, y para el análisis de la reserva de ley en el sector minero. \\
\hline

4 &
Guía de Jurisprudencia Constitucional sobre Derechos de la Naturaleza. Centro de Documentación de la Corte Constitucional del Ecuador (CEDEC). 2023 &
Documento de sistematización jurisprudencial. No es sentencia autónoma ni normativa vinculante &
Sistematiza diez sentencias constitucionales sobre derechos de la naturaleza, incluyendo referencias a decisiones cuyo texto completo no ha sido incorporado de manera independiente al corpus documental de este trabajo. &
Instrumento de consulta secundaria. Permite acceder al resumen de jurisprudencia complementaria sobre DDN. No reemplaza las sentencias individuales. \\
\hline

5 &
Comunicado institucional. Admisión de demandas y suspensión provisional de normas en leyes de reciente promulgación. Corte Constitucional del Ecuador. Agosto de 2025 &
Comunicado institucional contextual. No es sentencia. Refiere a leyes ajenas al sector minero (Ley de Integridad Pública, Ley de Inteligencia, Ley de Solidaridad Nacional) &
Ilustra el mecanismo procesal de admisión de acciones de inconstitucionalidad y suspensión provisional de normas ante la CCE. &
Fuente contextual exclusivamente. Ilustra el proceso constitucional aplicable, pero no tiene relación directa con la Ley de Fortalecimiento de Minería y Energía de 2026. No debe usarse como evidencia de acción constitucional sobre la reforma minera. \\
\hline

\end{longtable}
\endgroup

\noindent\textit{Nota.} Las sentencias No.\ 1149-19-JP/21, 32-17-IN/21 y 22-18-IN/21 son precedentes constitucionales vinculantes para todos los poderes del Estado, incluyendo la Asamblea Nacional en su función legislativa. La Guía CEDEC 2023 y el Comunicado institucional de agosto de 2025 son documentos de sistematización y comunicación institucional respectivamente; no tienen carácter de sentencia ni valor normativo autónomo. El Comunicado de agosto de 2025 refiere a leyes distintas al sector minero, razón por la cual no se utiliza como evidencia de decisión constitucional sobre la Ley Orgánica para el Fortalecimiento de los Sectores Estratégicos de Minería y Energía (2026).
